% !TEX program = xelatex
\documentclass[10pt,a4paper]{article}

% ---- Пакеты для XeLaTeX и кириллицы ----
\usepackage{fontspec}
\usepackage{polyglossia}
\setmainlanguage{russian}
\setotherlanguage{english}

% ---- Шрифт с кириллицей ----
\setmainfont{Liberation Serif}
\setsansfont{Liberation Sans}
\setmonofont{Liberation Mono}

% ---- Поля и геометрия ----
\usepackage[top=1.5cm, bottom=1.5cm, left=1.5cm, right=1.5cm]{geometry}
\usepackage{parskip}
\setlength{\parindent}{0pt}

% ---- Продвинутые ссылки ----
\usepackage{hyperref}
\hypersetup{
    colorlinks=true,
    linkcolor=blue,
    urlcolor=blue,
    citecolor=blue
}

% ---- Полезные пакеты ----
\usepackage{enumitem}
\setlist{nosep, left=0pt}
\usepackage{tabularx}
\usepackage{array}
\usepackage{booktabs}
\usepackage{xcolor}
\usepackage{titlesec}

% ---- Настройка заголовков разделов ----
\titleformat{\section}{\Large\bfseries\uppercase}{}{0em}{}[\vspace{-0.3em}\rule{\textwidth}{0.5pt}]
\titlespacing*{\section}{0pt}{0.8em}{0.3em}

\titleformat{\subsection}{\large\bfseries}{}{0em}{}
\titlespacing*{\subsection}{0pt}{0.4em}{0.2em}

% ---- Пользовательские команды ----
\newcommand{\resumeItem}[2]{\textbf{#1} --- #2\\}
\newcommand{\projectEntry}[3]{\textbf{#1} \hfill \textit{#2}\par #3\par\vspace{0.3em}}

% ---- Общие данные (имя, телефон, email и т.д.) ----
% ============================================================
% ОБЩИЕ ПЕРСОНАЛЬНЫЕ ДАННЫЕ
% ============================================================
% Меняешь здесь — меняется во всех резюме

\newcommand{\PersonName}{Кузык Иван Романович}
\newcommand{\PersonPhone}{+7 (916) 574-52-99}
\newcommand{\PersonEmail}{vanya.kuzyk1@gmail.com}
\newcommand{\PersonGitHub}{github.com/vanya890}
\newcommand{\PersonCity}{г.~Москва, м.~Бульвар Дмитрия Донского}
\newcommand{\PersonCitizenship}{РФ}
\newcommand{\PersonDrivingLicense}{Категория B}

% ============================================================
% КОМАНДА ДЛЯ ЗАГОЛОВКА РЕЗЮМЕ
% ============================================================
% Использование: \makeHeader{Должность}
\newcommand{\makeHeader}[1]{%
\begin{center}
    {\Huge\bfseries \PersonName}\\[0.3em]
    {\large #1}\\[0.5em]
    \begin{tabular}{ll}
        Телефон: & \PersonPhone \\
        Email:   & \href{mailto:\PersonEmail}{\PersonEmail} \\
        GitHub:  & \href{https://\PersonGitHub}{\PersonGitHub} \\
        Город:   & \PersonCity \\
        Гражданство: & \PersonCitizenship
    \end{tabular}
\end{center}
}

% ============================================================
% ОБЩЕЕ ОБРАЗОВАНИЕ
% ============================================================
% Использование: \makeEducation
\newcommand{\makeEducation}{%
\section{Образование}

\projectEntry{Московский политехнический университет --- Факультет информационных технологий}{2022--2026}
{Бакалавриат, направление «Информатика и вычислительная техника», специализация «Киберфизические системы». ВКР «Разработка образовательного 3D-симулятора робототехники» --- отлично. Средний балл --- 5.0.}

\projectEntry{ГБПОУ «Образовательный комплекс «Юго-Запад»}{2018--2022}
{Среднее специальное образование, специальность «Информационные системы и программирование». Диплом с отличием. Демонстрационный экзамен по компетенции «Разработка компьютерных игр и мультимедийных приложений» --- отлично.}
}

% ============================================================
% ОБЩАЯ ДОПОЛНИТЕЛЬНАЯ ИНФОРМАЦИЯ
% ============================================================
% Использование: \makeFooter{Ключевые компетенции}{Личные качества}
\newcommand{\makeFooter}[2]{%
\section{Дополнительная информация}

\begin{itemize}
    \item \textbf{GitHub:} \href{https://\PersonGitHub}{\PersonGitHub} --- #1
    \item \textbf{Ключевые компетенции:} #2
    \item \textbf{Личные качества:} Обучаемость, ответственность, умение работать с большим объёмом информации, системное мышление, ориентация на результат
    \item \textbf{Водительские права:} \PersonDrivingLicense
\end{itemize}
}

\begin{document}

% ============================================================
% ЗАГОЛОВОК
% ============================================================
\makeHeader{Разработчик веб- и мультимедийных приложений (Fullstack)}

% ============================================================
% О СЕБЕ
% ============================================================
\section{О себе}

Fullstack-разработчик с опытом создания веб-приложений, desktop-программ и AI-сервисов. Выпускник Московского политехнического университета по направлению «Информатика и вычислительная техника» (киберфизические системы). Владею современным стеком: \textbf{React, TypeScript, Node.js, PHP (Laravel), Python, Docker}. Разработал и опубликовал более 10 проектов на GitHub, включая образовательный 3D-симулятор робототехники (ВКР), AI-чат с персонажами и AI-визуальную новеллу. Ориентирован на результат, быстро осваиваю новые технологии, умею работать как в команде, так и самостоятельно.

% ============================================================
% ОПЫТ РАБОТЫ (ПРОЕКТНАЯ ДЕЯТЕЛЬНОСТЬ)
% ============================================================
\section{Опыт работы}

\subsection{Проектная деятельность / Фриланс}

\projectEntry{Fullstack-разработчик (EduRoboSim — ВКР)}{2025--2026}
{Разработал образовательный 3D-симулятор робототехники на стеке \textbf{Tauri (Rust + WebView2) + React 18 + TypeScript + Three.js (R3F)}. Спроектировал и реализовал собственный язык программирования (EduScript) с лексером, парсером (AST) и виртуальной машиной на TypeScript. Внедрил физический движок 2D OBB-коллизий, систему датчиков (8 дальномеров, энкодер, гироскоп) и редактор уровней в формате YAML. Применены принципы SOLID и паттерн MVC.}

\projectEntry{Fullstack-разработчик (AI Chat Game)}{2024--2025}
{Создал веб-приложение для чата с AI-персонажами на \textbf{React 18 + TypeScript + Vite}. Реализовал систему памяти (кратко- и долгосрочную), визуальный редактор персонажей, дашборд аналитики токенов. Интегрировал поддержку нескольких AI-провайдеров: OpenAI, Ollama, LM Studio, Groq. Использовал Zustand для управления состоянием, IndexedDB для хранения данных, i18next для мультиязычности.}

\projectEntry{Fullstack-разработчик (InfinityAI Novel)}{2025--2026}
{Разработал AI-визуальную новеллу на \textbf{React 19 + TypeScript + Vite 7 + Tailwind CSS 4}. Интегрировал генерацию сюжета через OpenRouter API и локальные модели LM Studio с JSON Schema structured output. Реализовал генерацию портретов персонажей через ComfyUI, систему автосохранения и серверное хранение данных.}

\projectEntry{Разработчик (laravelShop --- диплом колледжа)}{2022}
{Разработал интернет-магазин на \textbf{Laravel (PHP)} с каталогом товаров, корзиной, оформлением заказов и административной панелью. Использовал MySQL, Blade-шаблоны, Bootstrap.}

\projectEntry{Разработчик (Minecraft Server Project)}{2024--2025}
{Развернул полностью контейнеризованный игровой сервер на \textbf{Purpur 1.21} с лендинг-сайтом (\textbf{React + TypeScript + Vite}) в Docker-окружении. Разработал Windows-лаунчер на Electron с автоустановкой Java и параллельной загрузкой ассетов.}

\subsection{Практика}

\projectEntry{Техник-программист (ГБПОУ «Юго-Запад»)}{Апрель--Июнь 2022}
{Проверка работоспособности ПК, обновление и установка ОС и ПО, настройка сетевых подключений, установка принтеров и периферийного оборудования, помощь в организации ЕГЭ.}

\projectEntry{Оператор БД (Издательство «Просвещение»)}{Май--Июль 2021}
{Заполнение базы данных информацией, описание изображений для поиска по ключевым словам, подготовка данных для редакций.}

\projectEntry{Оператор ПК (ГБПОУ «Юго-Запад»)}{Ноябрь 2020 --- Январь 2021}
{Ввод и редактирование данных в программе \textbf{1С:Библиотека}, обучение сотрудников работе с системой на базовом уровне.}

% ============================================================
% ОБРАЗОВАНИЕ
% ============================================================
\makeEducation

% ============================================================
% НАВЫКИ
% ============================================================
\section{Навыки}

\begin{tabular}{p{3.5cm} p{12cm}}
    \textbf{Языки программирования:} & TypeScript, JavaScript, PHP, Python, C\#, Rust (базовый), SQL \\
    \textbf{Frontend:} & React 18/19, TypeScript, Vite, Tailwind CSS, Three.js / React Three Fiber, A-Frame, HTML5, CSS3 \\
    \textbf{Backend:} & Node.js, Express.js, Laravel, REST API, WebSocket \\
    \textbf{Базы данных:} & MySQL, SQLite, IndexedDB, Oracle (базовый) \\
    \textbf{Инструменты:} & Git, Docker, Docker Compose, Nginx, CI/CD \\
    \textbf{Фреймворки:} & Tauri (Rust + WebView2), .NET (Windows Forms), Unity \\
    \textbf{AI/ML:} & OpenAI API, Ollama, LM Studio, ComfyUI, LangChain (базовый) \\
    \textbf{Методологии:} & Agile, Scrum, SOLID, MVC, ООП \\
    \textbf{Языки:} & Русский --- родной, Английский --- A2 (элементарный)
\end{tabular}

% ============================================================
% ДОПОЛНИТЕЛЬНАЯ ИНФОРМАЦИЯ
% ============================================================
\makeFooter{более 16 репозиториев с открытым исходным кодом}{Разработка веб-приложений (SPA, SSR), проектирование архитектуры ПО, создание REST API, работа с базами данных, контейнеризация (Docker), интеграция AI-моделей}

\end{document}